\chapter{RESEARCH METHODOLOGY}
\label{ChapterResearchMethodology}

This chapter describes the research methodology adopted for the entire project, including data collection, feature engineering, model architecture design, training procedures, and evaluation framework.

\section{Research Design Overview}

The project follows an experimental research design with the following phases:

\begin{enumerate}
    \item \textbf{Data Collection:} Acquisition and preprocessing of high-frequency intraday data from the NSE.
    \item \textbf{Feature Engineering:} Construction of 47 features from raw OHLCV data organized into six categories.
    \item \textbf{Regime Detection:} Application of Hidden Markov Models for market state identification.
    \item \textbf{Model Development:} Implementation of the TFT architecture with two ablation variants.
    \item \textbf{Baseline Benchmarking:} Implementation of HAR-RV and GARCH family baselines.
    \item \textbf{Evaluation:} Comprehensive performance assessment using multiple metrics across regime conditions.
\end{enumerate}

\section{Dataset Description}

The dataset comprises 5-minute Open-High-Low-Close-Volume (OHLCV) data from fourteen highly liquid equities listed on the National Stock Exchange of India. The chosen securities represent multiple sectors to ensure diverse volatility dynamics:

\begin{itemize}
    \item \textbf{Banking:} HDFCBANK, ICICIBANK, AXISBANK, SBIN
    \item \textbf{Information Technology:} INFY, TCS, WIPRO
    \item \textbf{Energy:} RELIANCE, COALINDIA
    \item \textbf{Infrastructure:} LT, ADANIPORTS
    \item \textbf{Defence:} HAL
    \item \textbf{Metals:} TATASTEEL
    \item \textbf{Consumer:} TITAN
\end{itemize}

The dataset spans approximately ten years, giving $\sim$194,000 bars per stock and approximately 2.65 million rows in total. Indian equity trading runs from 9:15~AM to 3:30~PM IST, yielding 75 five-minute bars per trading day.

\subsection{Data Preprocessing}

The raw data underwent several preprocessing steps:

\begin{enumerate}
    \item \textbf{Outlier Removal:} Price bars with returns exceeding 10 standard deviations from the rolling mean were flagged and corrected using linear interpolation.
    \item \textbf{Missing Value Imputation:} Gaps in the time series (due to trading halts, data feed issues) were filled using forward-fill for prices and zero-fill for volumes.
    \item \textbf{Intraday Seasonality Adjustment:} The well-documented U-shaped intraday volatility pattern (high at open, low mid-day, rising toward close) was accounted for through normalization.
    \item \textbf{Train-Validation Split:} A strictly chronological split was enforced with no random shuffling, ensuring temporal integrity.
\end{enumerate}

\section{Feature Engineering Pipeline}

A total of 47 features were engineered from the raw price and volume data, organized into six categories as described below.

\subsection{Multi-Scale Realized Volatility}

Realized volatility is computed as the square root of the sum of squared high-frequency intraday returns $(r_{t,i})$:

\begin{equation}
RV_h = \sqrt{\sum_{i=1}^{h} r_{t,i}^2} \times \sqrt{252}
\end{equation}

This is annualized and calculated over multiple horizons: $h \in \{12, 38, 75, 375\}$ bars, corresponding to approximately 1 hour, half-day, 1 day, and 1 week respectively. This multi-scale approach captures volatility dynamics at different temporal resolutions.

\subsection{HAR-RV Lag Features}

Following the framework of Corsi \cite{corsi2009}, we explicitly supply RV lags at offsets corresponding to:
\begin{itemize}
    \item 1 bar (5 minutes)
    \item 12 bars (1 hour)
    \item 75 bars (1 day)
    \item 375 bars (1 week)
\end{itemize}

These lag features capture the heterogeneous nature of volatility memory, where different market participants operate at different frequencies.

\subsection{Jump Detection}

Based on the framework of Barndorff-Nielsen and Shephard \cite{bns2004}, we isolate continuous and jump components using Bipower Variation (BPV):

\begin{equation}
BPV_t = \frac{\pi}{2} \sum_{i=2}^{75} |r_{t,i}| \cdot |r_{t,i-1}|
\end{equation}

The jump component is then extracted as:
\begin{equation}
J_t = \max(RV_{75}^2 - BPV_t,\; 0)
\end{equation}

Jump features help the model distinguish between smooth volatility evolution and sudden discontinuities caused by news events, earnings releases, or macroeconomic announcements.

\subsection{Range-Based Estimators}

Three efficient variance estimators exploiting the full price range information are computed:

\textbf{Parkinson Estimator} \cite{parkinson1980}:
\begin{equation}
\sigma_{PK}^2 = \frac{1}{4\ln(2)} \left( \ln\frac{H}{L} \right)^2
\end{equation}

\textbf{Garman-Klass Estimator} \cite{garman1980}:
\begin{equation}
\sigma_{GK}^2 = \frac{1}{2} \left( \ln\frac{H}{L} \right)^2 - (2\ln 2 - 1) \left( \ln\frac{C}{O} \right)^2
\end{equation}

\textbf{Rogers-Satchell Estimator} \cite{rogers1991}:
\begin{equation}
\sigma_{RS}^2 = \ln\!\left(\frac{H}{C}\right) \ln\!\left(\frac{H}{O}\right) + \ln\!\left(\frac{L}{C}\right) \ln\!\left(\frac{L}{O}\right)
\end{equation}

These estimators provide more efficient variance estimates than close-to-close returns by incorporating information from daily High ($H$), Low ($L$), Open ($O$), and Close ($C$) prices.

\subsection{GARCH and Cross-Asset Features}

A standard GARCH(1,1) model \cite{bollerslev1986} is fitted per stock to extract conditional variance:

\begin{equation}
\sigma_t^2 = \omega + \alpha\, \epsilon_{t-1}^2 + \beta\, \sigma_{t-1}^2
\end{equation}

Additionally, a DCC-inspired proxy captures time-varying correlations between each stock and the market portfolio. This cross-asset feature enables the model to detect contagion effects and market-wide volatility spillovers.

\subsection{Temporal Encodings}

Cyclical sine/cosine transformations encode the time of day and day of week, alongside binary session flags:
\begin{itemize}
    \item \texttt{is\_opening}: First 15 minutes of the trading session
    \item \texttt{is\_closing}: Last 15 minutes of the trading session
\end{itemize}

These features capture the well-documented intraday volatility patterns in equity markets.

\section{Target Variable Construction}

The prediction target is next-day realized volatility computed over the next 75 five-minute bars:

\begin{equation}
y_t = \log\!\left(1 + \sqrt{\sum_{i=t+1}^{t+75} r_i^2} \times \sqrt{252}\right)
\end{equation}

The log-transformation stabilizes the variance of the target and ensures non-negativity. By construction, this target uses exclusively future returns, ensuring zero temporal overlap between features and the prediction target.

\section{Regime Detection Methodology}

A four-state Gaussian Hidden Markov Model is estimated on the market-average log-returns and the daily realized volatility. The HMM parameters---transition probabilities and emission distributions---are learned using the Baum-Welch algorithm (Expectation-Maximization). The emission means of the inferred hidden states are then ordered to define four volatility regimes:

\begin{enumerate}
    \item \textbf{Low Volatility:} Calm market conditions with minimal price fluctuations.
    \item \textbf{Normal Volatility:} Typical market conditions with moderate price movements.
    \item \textbf{High Volatility:} Sustained elevated volatility, often during trending markets.
    \item \textbf{Extreme Volatility:} Crisis-level volatility during market dislocations.
\end{enumerate}

These regime labels are fed as categorical features into the TFT, allowing the model to condition its predictions on the prevailing market state.

\section{Evaluation Framework}

The following metrics are used for comprehensive evaluation:

\begin{itemize}
    \item \textbf{$R^2$ (Coefficient of Determination):} Measures the proportion of variance explained.
    \item \textbf{RMSE (Root Mean Square Error):} Measures average prediction error magnitude.
    \item \textbf{MAPE (Mean Absolute Percentage Error):} Scale-independent error measure.
    \item \textbf{DA (Directional Accuracy):} Proportion of correct volatility direction predictions.
    \item \textbf{QLIKE:} Quasi-likelihood loss function specific to volatility forecasting; lower values indicate better predictive density.
\end{itemize}
